\documentclass[../../main.tex]{subfiles}
\begin{document}

\chapter{Code refactoring, maintainability, and technical debt}
\label{ch:code_refactoring_maintainability}

\section{Chapter overview and learning outcomes}
Refactoring improves the internal quality and maintainability of code without altering its external observable behavior. This chapter investigates how Large Language Models detect architectural code smells, perform semantic-preserving transformations, and reduce technical debt while respecting formal maintainability metrics.

Upon completing this chapter, you will be able to:
\begin{itemize}
    \item Identify classic Fowler code smells (God Class, Long Method, Feature Envy, Shotgun Surgery).
    \item Compute formal software metrics: McCabe's Cyclomatic Complexity, Cognitive Complexity, and the Maintainability Index (MI).
    \item Formulate prompt chains for automated refactoring transformations (Extract Method, Introduce Parameter Object).
    \item Establish semantic equivalence verification via regression suites and differential testing.
    \item Quantify technical debt reduction achieved by AI-assisted code transformations.
\end{itemize}

\section{Software complexity metrics}
Before and after applying an AI refactoring step, engineers must verify that structural complexity metrics decrease.

\begin{definition}{McCabe cyclomatic complexity}{cyclomatic_complexity}
Given the Control Flow Graph (CFG) of a program with $E$ directed edges, $N$ nodes, and $P$ connected components, McCabe's Cyclomatic Complexity $M$ is defined as:
\begin{equation}
\label{eq:cyclomatic}
M = E - N + 2 P.
\end{equation}
For a single function ($P=1$), $M$ equals the number of linearly independent paths through the code.
\end{definition}

\begin{definition}{Maintainability index (MI)}{maintainability_index}
The classic Maintainability Index combines Halstead Volume $V$, Cyclomatic Complexity $G$, lines of code $\text{LOC}$, and comment ratio $\text{CM}$:
\begin{equation}
\text{MI} = 171 - 5.2 \ln V - 0.23 G - 16.2 \ln(\text{LOC}) + 50 \sin\left(\sqrt{2.4 \cdot \text{CM}}\right).
\end{equation}
Values above $85$ denote high maintainability; values below $65$ signify severe maintenance risks.
\end{definition}

\begin{examinsight}[Exam focus: behavioral equivalence proofs]
A standard exam question queries: \emph{"How can we guarantee that an LLM-refactored method has not altered program behavior?"}
Unlike formal refactoring engines with AST rewriting rules, LLMs are stochastic and can silently modify operator precedence or variable state. To prove equivalence:
\begin{enumerate}
    \item Run mutation-tested regression test suites before and after refactoring.
    \item Perform differential fuzzing on randomized inputs.
    \item When feasible, employ bounded model checking or SMT solvers (e.g., Z3) to verify that outputs coincide for all symbolic inputs.
\end{enumerate}
\end{examinsight}

\end{document}
